% Korean edition of sections/s2_5_orders.tex.
\seclabel
\claim{\secnum{2.5}{sec:orders}\quad 주문은 REST로 넣고 피드에서 확인하며, 크기는 \code{wei}가 정한다}

\lead{서버는 개장 전에 넣은 주문도 받는다.} 08:30:05에 웹 클라이언트 코드에서 가져온 본문으로 \code{AddOrder} 요청을 보냈다. 본문에는 날짜(\code{d}), 포트폴리오(\code{pfId}), 종목코드(\code{prodId} = \code{A079550}, LIG넥스원), 방향(\code{ls} = L), 진입/청산 코드(\code{ex} = E는 포지션을 열거나 늘림, X는 줄임, \ref{sec:wrappertests}절), 주문 크기 필드(\code{wei} = 0.3)가 들어 있다. 주문이 어떻게 거래될지는 두 묶음의 필드가 더 정한다. 가격 유형 필드는 고정 가격 지정가(\code{limitPrc}), 시장가에서 $n$호가 떨어진 지정가(\code{limitIdx}), 스톱 가격(\code{stopPrc}) 중 하나를 고른다. 이 보고서의 다른 곳에서 ‘정지’는 진행 중인 주문을 취소하는 것을 뜻하며, 웹 페이지가 \code{CancelOrderBatch}를 부르는 말이다. 스톱 가격과는 관계없다. 나머지 두 필드인 시간 필드 \code{hm0}, \code{hm1}은 주문을 즉시 보내거나 체결을 시간대에 나눠 분산한다. 이 주문은 09:00\textasciitilde 09:10이었다. 서버는 HTTP 200과 \texttt{\{"data":"","warnings":[]\}}로 답했다.

\lead{응답에 주문 ID가 없으므로 래퍼는 실시간 피드에서 자기 주문을 찾는다.} 주문 직후의 피드 응답은 인자가 \code{null}이었고(\ref{sec:feed}절), 08:33까지 피드의 주문 목록(\code{ords})에 ID 1649882인 새 레코드가 생겼다. 래퍼는 쓰기 전에는 없었고 \code{place()}가 아직 돌려주지 않은 ID 가운데서 종목코드, 방향, 크기가 맞는 새 레코드를 찾으며, 쓰기와 검색 내내 잠금 하나를 잡는다. 덕분에 이 프로세스의 같은 주문 2건이 같은 ID에 묶이지 않는다. 첫 주문이 REST로만 확인된 경우에도 마찬가지다. 맞는 새 레코드가 둘 이상이면 \code{place()}는 하나를 고르지 않고 \code{AmbiguousState}를 던진다. 같은 순간 웹 페이지나 다른 프로세스에서 넣은 똑같은 주문 1건은 응답에 ID가 없으므로 구분할 수 없다. 그래서 래퍼는 실행 중에 계정당 주문 제출자가 하나라고 가정한다.

\lead{응답을 잃으면 래퍼는 주문이 있는지 확신할 수 없으므로, 서버가 받았을 수 있는 주문은 다시 보내지 않는다(\figref{orderstates}).} 타임아웃, 5xx, 또는 요청이 나간 뒤 연결이 끊기면, 래퍼는 15\,s 동안 피드에서 주문을 찾고, 다음으로 REST \code{Portfolio/Orders}에서 쓰기 5\,s 전 이후에 생성된 레코드(시계 오차 여유)를 찾는다. 둘 다 없으면 \code{AmbiguousState}를 던지고 확인을 호출자에게 맡긴다. 다시 보내면 주문이 두 번 들어갈 수 있고, 추가 주문은 래퍼가 허용하는 하루 80건 중 1건을 쓴다. 래퍼가 \code{AddOrder}를 다시 보내는 경우는 두 가지뿐이다. 첫째, 연결이 아예 맺어지지 않았다면(\code{ConnectError}, \code{ConnectTimeout}, \code{PoolTimeout}) 요청이 서버에 닿았을 수 없으므로 한 번 더 보낸다. 그 재전송도 실패하면 같은 피드·REST 검색을 거치며 세 번째 전송은 없다. 둘째, 401을 받으면 토큰을 갱신하고 한 번 더 보낸다. ASP.NET은 주문 처리기가 실행되기 전에 요청을 거부하기 때문이다(\ref{sec:auth}절). 400은 절대 재시도하지 않으며, 어떤 요청에서든 403이나 429를 받으면 이후 모든 쓰기를 멈춘다(\ref{sec:design}장).

% Korean edition of figs/orderstates.tex.
\fig[wide]{orderstates}{서버가 받았을 수 있는 주문은 다시 보내지 않는다. 서버의 주문은 체결이든 정지든 \code{Done}으로 끝난다}[위: \code{AddOrder} 호출 하나가 응답 종류별로 하는 일. 아래: 서버에서 주문의 \code{state}(이름은 웹 클라이언트 코드에서)]{%
\begin{tikzpicture}[x=1mm,y=1mm,
  f/.style={font=\fontsize{7}{8.4}\selectfont},
  st/.style={rectangle,rounded corners=1.6mm,draw=navy,line width=0.6pt,f,text=navy,align=center,minimum width=24mm,minimum height=8.5mm,inner sep=1pt,fill=white},
  sv/.style={st,rounded corners=4.2mm,minimum width=22mm},
  un/.style={sv,draw=greyD,dashed,text=greyD},
  term/.style={st,fill=navy,text=white},
  tr/.style={-{Stealth[length=1.6mm]},line width=0.6pt,navy},
  trU/.style={-{Stealth[length=1.6mm]},line width=0.5pt,greyD,dashed},
  lb/.style={f,text=navy,inner sep=1pt,align=center},
  lbU/.style={f,text=greyD,inner sep=1pt,align=center},
  band/.style={font=\fontsize{7.5}{9}\selectfont\bfseries,text=navy,anchor=west}]
% ------------------------------------------------ wrapper band
\node[band] at (-22,33) {래퍼: \code{AddOrder} 호출 하나};
\node[st] (S) at (32,0) {\code{AddOrder}\\보내는 중};
\node[st] (F) at (84,0) {피드에서\\주문 찾는 중};
\node[st] (Q) at (134,0) {REST \code{Orders}\\조회 중};
\node[term] (T4) at (32,25) {\code{RequestError}\\발생};
\node[term] (TA) at (134,25) {\code{AmbiguousState}:\\주문 재전송 안 함};
\node[term] (TI) at (116,-17) {주문 ID 반환};
% self loops on S
\draw[tr] ([yshift=2mm]S.west) to[out=160,in=110,looseness=6] ([xshift=-7mm]S.north);
\node[lb,anchor=east,align=right] at (13,9) {401: 토큰 갱신 후\\한 번 더 보냄};
\draw[tr] ([yshift=-2mm]S.west) to[out=-160,in=-110,looseness=6] ([xshift=-7mm]S.south);
\node[lb,anchor=east,align=right] at (13,-9) {연결이 맺어지지 않음:\\한 번 더 보냄. 그것도\\실패하면 주문을 찾음};
% S -> T4, S -> F
\draw[tr] (S.north) -- node[lb,right]{400: 재시도 안 함} (T4.south);
\draw[tr] (S.east) -- node[lb,above]{주문 ID 없는\\200} node[lb,below,align=center]{또는 응답 유실:\\타임아웃, 5xx,\\보낸 뒤\\연결 끊김} (F.west);
% F -> Q, F -> TI, Q -> TA, Q -> TI
\draw[tr] (F.east) -- node[lb,above]{15\,s 안에} node[lb,below]{못 찾음} (Q.west);
\draw[tr] (Q.north) -- node[lb,right,align=left]{쓰기 이후 레코드에\\없음} (TA.south);
\draw[tr] (F.south) |- (TI.west);
\node[lb,anchor=north] at (92,-18.2) {찾음: 같은 종목코드,\\\code{ex}, 크기};
\draw[tr] (Q.south) |- node[lb,pos=0.25,right,xshift=0.8mm,align=left]{찾음} (TI.east);
% ------------------------------------------------ separator
\draw[greyM,line width=0.4pt] (-22,-27) -- (158,-27);
\node[band] at (-22,-31) {서버: 주문의 \code{state}};
% ------------------------------------------------ server band
\node[un] (G) at (-10,-44) {\code{Generated}};
\node[sv] (A) at (28,-44) {\code{Accepted}};
\node[sv] (W) at (70,-44) {\code{Working}};
\node[sv] (D) at (112,-44) {\code{Done}};
\node[un] (X) at (9,-62) {\code{Deleted}};
\draw[trU] (G) -- (A);
\draw[trU] (A) -- node[lbU,above,yshift=0.8mm]{관찰 못함} (W);
\draw[trU] (G) -- (X);
\draw[trU] (A) -- (X);
\node[lbU,anchor=east,align=right] at (-3,-58) {\code{DeleteOrder}};
\draw[tr] (W.north east) to[out=25,in=155] node[lb,above]{마지막 주식 체결} (D.north west);
\draw[tr] (W.east) -- node[lb,below,align=center]{정지} (D.west);
\draw[tr] (A.south) to[out=-35,in=-145] node[lb,below,align=center]{즉시 정지(주문 1650719)} (D.south);
\node[lb,anchor=west,align=left] at (125,-44) {\code{cnclT} null이면 체결,\\\code{cnclT} 설정이면 정지,\\거기에 \code{cumQty} $>$ 0이면\\일부 체결 후 정지};
% ------------------------------------------------ legend
\begin{scope}[shift={(-22,-78)}]
\node[st,minimum width=8mm,minimum height=4.5mm] at (4,0) {}; \node[f,anchor=west] at (9,0) {래퍼 단계};
\node[term,minimum width=8mm,minimum height=4.5mm] at (32,0) {}; \node[f,anchor=west] at (37,0) {호출 종료};
\node[sv,minimum width=8mm,minimum height=4.5mm,rounded corners=2.2mm] at (63,0) {}; \node[f,anchor=west] at (68,0) {서버 상태, 관찰함};
\node[un,minimum width=8mm,minimum height=4.5mm,rounded corners=2.2mm] at (101,0) {}; \node[f,anchor=west] at (106,0) {클라이언트 코드에만 있음};
\draw[tr] (145,1.6) -- (152,1.6); \node[f,anchor=west] at (153,1.6) {규칙 또는 관찰};
\draw[trU] (145,-1.8) -- (152,-1.8); \node[f,anchor=west] at (153,-1.8) {관찰 못함};
\end{scope}
\end{tikzpicture}}{래퍼는 첫 주문이 서버에 닿았는지 알 수 없으면 주문을 다시 보내지 않는다. 한 번 더 보내는 경우는 중복을 만들 수 없는 두 가지뿐이다. 연결이 아예 맺어지지 않은 경우(\code{ConnectError}, \code{ConnectTimeout}, \code{PoolTimeout})와, ASP.NET이 요청이 핸들러에 닿기 전에 돌려주는 401이다(D36). 연결 실패 뒤 재전송도 실패하면 주문은 피드·REST 검색으로 넘어가며, 세 번째 전송은 없다. 어떤 요청에서든 403이나 429를 받으면 \code{BlockedError}를 던지고 모든 쓰기를 멈춘다(그림에 없음). 잠금 하나가 보내기와 찾기를 덮고 이미 돌려준 ID는 건너뛰므로, 같은 주문 2건이 같은 주문 ID에 대응될 수 없다. 서버 상태: \code{Accepted}는 주문 1650719를 넣은 직후, \code{Working}은 08:30 주문에서 보았다. \code{Generated}와 \code{Deleted}는 클라이언트 코드에만 나온다. 끝난 주문은 모두 \code{Done}으로 보이므로 \code{cnclT}와 \code{cumQty}로 체결과 정지를 구분한다.}{\code{pytimefolio}의 \code{orders}, \code{transport} 모듈; 2026년 11월 24일 실시간 피드와 \code{Orders} 응답; DART Quant Team}


\lead{주문의 \code{wei}는 순자산(\code{na}) 대비 주문 크기(\%)다. 목표 비중으로 읽었다면 주식 수가 달랐을 것이다.} 주문 전 포트폴리오는 LIG넥스원 3주(순자산의 0.13\%)를 갖고 있었고, 1주는 약 0.042\%였다. \code{wei} = 0.3을 목표로 읽으면 약 4주를 더 사라는 뜻이고, 크기로 읽으면 7.1주어치다. 주문은 \code{qty2ord} = 7을 보고했으므로 래퍼는 목표를 주문으로 바꿀 때 목표 비중에서 현재 비중을 뺀다. 래퍼는 주문을 만들 때 \code{wei}를 0.01\% 단위로 한 번 반올림하고, 같은 값을 보내고 같은 값으로 대조한다. 서버는 주문 크기를 보낸 그대로 저장했다(0.04, 0.1, 0.3을 확인). 서버는 \code{qty2ord}를 정수 주식 수로 내림한다. 내림과 반올림을 구분하는 테스트는 \ref{sec:oneshare}절에 있다.

\lead{주문은 개장 전 \code{Working} 상태였고, 약 2분마다 1\textasciitilde 2주씩 체결되어 7분 37초 동안 7주가 체결됐다(\figref{fills}).} 3\,s마다 폴링한 피드에서 체결 수량은 09:00:30에 1주, 09:02:09에 3주, 09:04:10에 4주, 09:06:10에 6주, 09:08:07에 7주였고, 이때 상태가 \code{Done}으로 바뀌었다. 주문 시간대가 끝나기 1분 53초 전이다. 포지션은 10주, 포트폴리오의 0.43\%가 됐다.

% Korean edition of figs/fills.tex.
\fig{fills}{7주가 약 2분마다 1\textasciitilde 2주씩 체결되어 09:08:07에 끝났다}[주문 1649882(LIG넥스원, \code{wei} 0.3)의 체결 주식 수(\code{cumQty}), 2026년 11월 24일 실시간 피드에서 읽음]{%
\begin{tikzpicture}[x=0.15mm,y=4.6mm,font=\fontsize{7}{8.4}\selectfont]
% x: seconds after 09:00 (0.15 mm per s); y: shares filled
\fill[navyXL] (0,0) rectangle (600,7.6);
% grid and axes
\foreach \y in {1,...,7} \draw[greyM,line width=0.2pt] (-60,\y) -- (660,\y);
\draw[greyD,line width=0.4pt] (-60,0) -- (660,0);
\draw[greyD,line width=0.4pt] (-60,0) -- (-60,7.6);
\foreach \y in {0,...,7} \node[anchor=east,text=greyD] at (-60,\y) {\y};
\foreach \m/\l in {0/09:00,120/09:02,240/09:04,360/09:06,480/09:08,600/09:10}
  {\draw[greyD] (\m,0) -- (\m,-0.18); \node[anchor=north,text=greyD] at (\m,-0.2) {\l};}
\node[anchor=north,text=greyD] at (300,-1.05) {시각 (KST)};
\node[rotate=90,anchor=south,text=greyD] at (-128,3.5) {7주 중 체결 주식 수};
\node[text=navy,align=center,fill=navyXL,inner sep=1pt] at (470,1.5) {주문 시간대\\09:00\textasciitilde 09:10};
% step line
\draw[navy,line width=1pt] (-60,0) -- (30,0) -- (30,1) -- (129,1) -- (129,3) -- (250,3) -- (250,4) -- (370,4) -- (370,6) -- (487,6) -- (487,7) -- (660,7);
% fill times
\foreach \t/\n/\l in {30/1/09:00:30,129/3/09:02:09,250/4/09:04:10,370/6/09:06:10}
  {\fill[navy] (\t,\n) circle (0.75mm); \node[anchor=south east,text=navy,inner sep=1.5pt] at (\t,\n) {\l};}
\fill[navy] (487,7) circle (0.75mm);
\node[anchor=north west,text=navy,inner sep=2pt,font=\fontsize{7}{8.4}\selectfont\bfseries] at (487,7) {09:08:07 \code{Done}};
% position weight after each fill
\node[anchor=east,text=greyD] at (-60,8.35) {포지션 비중(\%)};
\foreach \t/\w in {30/0.17,129/0.26,250/0.30,370/0.38,487/0.43}
  {\node[text=greyD] at (\t,8.35) {\w};}
\end{tikzpicture}}{점은 실시간 피드에서 읽은 체결 시각이다. 피드는 3\,s마다 폴링했으므로 각 시각은 $\pm$3\,s 정확도다. 음영은 주문 시간대(\code{hm0} 09:00\textasciitilde\code{hm1} 09:10)다. 주문은 08:33부터 \code{Working}이었고 마지막 주식이 체결되며 \code{Done}으로 바뀌었다. 첫 체결 후 7분 37초다. 맨 위 줄은 체결 후 포지션 비중(순자산 대비 \%)이며, 주문 전에 보유하던 3주를 포함한다. 평균 체결가 352,714.29원.}{실시간 피드 모니터 로그, DART Quant Team}


\lead{포지션 평균 단가에는 수수료 0.1\%가 들어 있다.} 주문 전 포지션은 평균 381,381원에 3주였고, 7주가 평균 352,714.29원에 체결됐다. 최종 평균 단가 361,561.2원은 (3\,×\,381,381 + 7\,×\,352,714.29\,×\,1.001)\,/\,10과 정확히 같으므로, \code{avgPrc}는 체결가에 수수료를 더한 값이다. 래퍼는 서버가 보고한 \code{avgPrc}로 손익을 계산하며 수수료를 다시 더하지 않는다.

\lead{주문 넣기, 정지, 조회는 실시간으로 기록했고, 나머지 주문 경로 3개는 클라이언트 코드에서만 읽었다(\figref{orderapi}).} 09:16\textasciitilde 09:22에 손으로 주문을 넣고 정지하는 동안 브라우저 개발자 도구로 페이지 자체의 요청을 기록했고, 나머지 본문은 JavaScript 번들에서 읽었다. \code{AddOrder}, \code{CancelOrderBatch}, \code{ValidateProduct}, \code{Orders}는 실제로 관찰했다. \code{CancelOrder}, \code{ModifyOrder}, \code{DeleteOrder}는 관찰하지 못했다. 실제 주문 대화상자는 \code{CancelOrderBatch}로 주문을 정지하고, 페이지는 미접수·실패 주문 목록에 있는 주문에만 정정과 삭제를 보여주므로, 이 3개 경로의 본문은 테스트되지 않았다. 래퍼는 \code{CancelOrderBatch}로 주문을 정지하며, ID를 쉼표로 잇고 이미 \code{Done}인 주문은 뺀다. 테스트되지 않은 3개 경로는 실서버에서 시험할 때까지 쓰지 않는다. 따라서 정정은 정지 후 새 주문으로 하며, 하루 한도에서 2건으로 센다.

% Korean edition of figs/orderapi.tex.
\fig[wide]{orderapi}{\code{AddOrder}, \code{CancelOrderBatch}, \code{Orders}는 실제로 캡처했고, 클라이언트 코드에서만 본 경로 3개는 시험하지 않았다}[\code{/api/Portfolio} 아래 경로(위)와 주문 본문의 가격·시간 필드(아래), 2026년 11월 24일 09:16\textasciitilde 09:22 KST 캡처]{%
{\tabfont\begin{tabularx}{\linewidth}{>{\RaggedRight\arraybackslash}p{26mm} >{\RaggedRight\arraybackslash}p{9mm} >{\RaggedRight\arraybackslash}X >{\RaggedRight\arraybackslash}p{40mm} >{\RaggedRight\arraybackslash}p{27mm}}
\thd{경로} & \thd{메서드} & \thd{매개변수 또는 본문} & \thd{응답} & \thd{출처}\\
\hl\code{AddOrder} & POST & \code{d}, \code{pfId}, \code{prodId}, \code{ls}, \code{ex}, \code{wei}, \code{exitAll}, \code{limitPrc}, \code{limitIdx}, \code{stopPrc}, \code{hm0}, \code{hm1} & 200, \texttt{\{"data":"","warnings":[]\}}. 주문 ID 없음 & 실제(매도 형식: \ref{sec:wrappertests}절)\\\rowsep
\hl\code{CancelOrderBatch} & POST & \code{d}, \code{pfId}, \code{ords}: 주문 ID를 쉼표로 이은 문자열 하나(예: \code{"1650197"}) & 200, 빈 본문. 페이지는 ‘정지’라고 표시 & 실제\\\rowsep
\hl\code{Orders} & GET & \code{d0}, \code{d1}, \code{pfId} & 200, 범위 안 모든 주문(\code{Done} 포함) & 실제\\\rowsep
\code{ValidateProduct} & GET & \code{d}, \code{prodId}, \code{entry=true} & 200, \texttt{\{"data":"","warnings":[]\}} & 실제\\\rowsep
\code{Session} & GET & \code{d}, \code{pfId}. 주문 대화상자가 열릴 때 호출 & 200, 시간대 09:00\textasciitilde 15:20 & 실제\\\rowsep
\code{CancelOrder} & POST & \code{d}, \code{pfId}, \code{ordId} & 관찰 못함 & 클라이언트 코드, 미시험\\\rowsep
\code{ModifyOrder} & POST & \code{d}, \code{ordId}, \code{wei}, \code{exitAll}, \code{limitPrc}, \code{limitIdx}, \code{stopPrc}, \code{hm0}, \code{hm1}. \code{pfId} 없음 & 관찰 못함 & 클라이언트 코드, 미시험\\\rowsep
\code{DeleteOrder} & POST & \code{d}, \code{pfId}, \code{ordId}. 미접수 주문에 제공 & 관찰 못함 & 클라이언트 코드, 미시험\\
\end{tabularx}}\par\vspace{5pt}
{\tabfont\begin{tabularx}{\linewidth}{>{\RaggedRight\arraybackslash}p{26mm} >{\RaggedRight\arraybackslash}p{52mm} >{\RaggedRight\arraybackslash}X}
\thd{필드} & \thd{값} & \thd{뜻}\\
\code{limitPrc} & 가격, \code{limitIdx} = null & 고정 가격 지정가 주문\\\rowsep
\code{limitIdx} & $n$ 또는 $-n$, $n$은 1\textasciitilde 10 & 상대 호가 쪽 $n$틱($n$) 또는 주문 자신의 호가 쪽($-n$)\\\rowsep
\code{stopPrc} & 가격 & 스톱 주문\\\rowsep
\code{hm0}, \code{hm1} & 현재 시각과 null, 또는 시작·종료 시각 & 즉시, 또는 \code{hm0}부터 \code{hm1}까지 시간대에 체결 분산\\
\end{tabularx}}}{강조 행: 래퍼가 주문을 넣고, 정지하고, 조회하는 데 쓰는 경로 3개. 래퍼는 주문마다 \code{ValidateProduct}를, 거래 시간대를 위해 \code{Session}을 호출하기도 한다. 위 표 아래쪽 경로 3개는 JavaScript 번들에만 나오며 래퍼는 호출하지 않는다. ‘실제’는 페이지에서 요청과 응답을 기록했다는 뜻이고, ‘클라이언트 코드’는 JavaScript 번들에서 본문을 읽었다는 뜻이다. 래퍼의 주문 호출은 \code{limitPrc}, \code{limitIdx}, \code{stopPrc} 중 정확히 하나를 받는다(\ref{sec:wrappertests}절).}{웹 클라이언트 코드, DART Quant Team}


\lead{끝난 주문은 모두 \code{Done}이 되므로, 어떻게 끝났는지는 \code{cnclT}와 \code{cumQty}로 읽어야 한다.} 번들은 상태 5개를 정의한다. \code{Generated}, \code{Accepted}, \code{Working}, \code{Done}, \code{Deleted}다. 페이지는 \code{state}가 \code{Done}이거나 정지 시각 \code{cnclT}가 null이 아니면 정정과 정지를 막는다. 09:31의 \code{Orders} 조회에서 정지한 삼성전자 주문 3건은 모두 \code{Done}이었고, \code{cnclT}가 설정돼 있었으며 4주 중 0주가 체결됐다(\code{qty2ord} 4 중 \code{cumQty} 0). 체결된 주문 2건과 같은 상태다. 그래서 래퍼는 \code{cnclT}가 설정된 \code{Done} 주문을 정지로, 0~$<$~\code{cumQty}이면 일부 체결 후 정지로 읽고, \code{cnclT}가 null인 \code{Done} 주문을 체결로 읽는다. \code{cnclT}가 null이면서 \code{cumQty}~$<$~\code{qty2ord}인 \code{Done} 주문은 아직 관찰하지 못했다. 시간대 주문이 체결되지 않은 채 \code{hm1}에 이르면 이렇게 될 수 있다. \code{wait()}은 이를 체결로 보고하며, 질문 Q5(\ref{sec:validity}절)에 올려 두었다. 관찰한 상태 순서는 두 가지다. 오전 주문들은 \code{Working}에서 \code{Done}으로, 10:33의 주문 1650719(\ref{sec:oneshare}절의 내림 테스트)는 \code{Accepted}에서 \code{Done}으로 갔다. 후자는 주문 직후 \code{Accepted}를 보였고 즉시 정지했다. \code{Generated}와 \code{Deleted}는 관찰하지 못했다.

\lead{이 절의 테스트에는 새 주문 5건과 정지 3건을 썼다.} 08:30 LIG넥스원 주문은 스크립트로 넣었고 체결됐다. 삼성전자 주문 2건은 손으로 넣었다. 09:13의 0.1\% 주문(\code{qty2ord} 4)은 정지했고, 09:15의 0.3\% 주문은 상대 호가 10틱 가격으로 넣어 13주 중 13주가 체결됐다. 요청을 기록하려고 체결되지 않도록 시장가보다 20\% 낮은 151,800원 지정가로 삼성전자 0.1\% 주문 2건(각 \code{qty2ord} 4)을 더 넣었다. 09:18(즉시)과 09:20(15:00\textasciitilde 15:20 시간대)에 넣었고 둘 다 1분 안에 정지했다. 그날 한도 대비 합계는 \ref{sec:wrappertests}절에서 센다.
\verdict{래퍼는 \code{AddOrder}마다 피드의 \code{ords}에서, 다음으로 \code{Portfolio/Orders}에서 새 레코드를 찾아 확인하며, 둘 다 없으면 다시 보내지 않고 \code{AmbiguousState}를 던진다. 연결이 맺어지지 않은 요청은 한 번만, 401로 거부된 요청은 한 번만 다시 보낸다(D11, D36). 주문의 \code{wei}는 목표에서 현재 비중을 뺀 값으로 보낸다(D13). \code{Done}은 끝난 것으로 읽는다. \code{cnclT}가 null이면 체결, 설정돼 있으면 정지, 그리고 0~$<$~\code{cumQty}이면 일부 체결 후 정지다(D15).}
